% ============================================================
%  Single section "Tasks and Rollouts": data construction + how agents are
%  run on the tasks.  No "benchmark", no leaderboard framing, no scoring in
%  the construction half.
%
%  Labels preserved / relocated:
%    sec:benchmark / sec:construction / sec:principles -> \section (aliases)
%    sec:taskconstruct -> "Tasks" paragraph
%    sec:stats         -> "Task types and composition" paragraph
%    sec:rollout       -> "Six-stage rollout" paragraph (now DEFINED here; the
%                         3 prior dangling \ref{sec:rollout} resolve to it)
%    sec:setup:agents / sec:setup:rollout -> "Agents" / "Rollout protocol"
%    sec:setup:eval    -> REMOVED from prose (was undefined); see note below
%    fig:stats         -> sits right after the composition paragraph

% ============================================================
\providecommand{\Set}{\mathbf{Set}}
\providecommand{\Lan}{\operatorname{Lan}}
\providecommand{\coel}[1]{\textstyle\int_{#1}}
\providecommand{\Kb}{\mathcal{K}}
\providecommand{\im}{\operatorname{im}}
\providecommand{\Prov}{\mathrm{Prov}}
\providecommand{\corpus}{\mathcal{C}}
\providecommand{\Tnov}{\mathcal{T}}
\providecommand{\Extract}{\mathcal{E}}
\providecommand{\Harn}{H}
\newtheorem{definition}{Definition}[section]

% ---------------- main text ----------------

\section{AutoResearch Tasks and Trajectories}
\label{sec:benchmark}
\label{sec:construction}  % legacy alias
\label{sec:principles}    % legacy alias; formalism now in Appendix~\ref{app:extraction}

\bench turns published papers into discovery tasks and runs agents on them end to end. From each
paper we build a task that states where the science stood and what was left unresolved---but no
method, so many paths through it are admissible---while the paper's published outcome is withheld. Tasks come in two types and carry domain and contribution labels, and other
fields extracted from the paper drive the filtering and authoring behind the scenes. Each task is
then run as a single autonomous rollout in a sandbox with code execution, and we log the complete
trajectory with produced artifacts as the unit of analysis. The rest of this section
details the tasks (\S~\ref{sec:taskconstruct}) and the rollouts (\S~\ref{sec:rollout}); extraction
prompts and the rollout environment are in Appendix~\ref{app:extraction}.

\newpage
\subsection{AutoResearch Tasks Mined from Venues}
\label{sec:taskconstruct}

\subsubsection{Task construction}
A paper is parsed into seven fields---\textsf{Premise}, \textsf{Tension}, \textsf{Motivation},
\textsf{Method}, \textsf{Experiment}, \textsf{KeyClaims}, \textsf{Conclusion}---and split into a
task instance $\tau_p=(q_p,\nu(p);\mathrm{target}_p)$. The query
$q_p=(\textsf{Premise},\textsf{Tension})$ states the prior literature and the anomaly left open by
it; the published outcome $\mathrm{target}_p=(\textsf{KeyClaims},\textsf{Conclusion})$ is withheld
from the query and kept as ground truth. The remaining five fields are not shown to the agent;
they drive construction: \textsf{Experiment} and \textsf{Method} gate whether a paper can become a
runnable task, \textsf{Motivation} and \textsf{Conclusion} determine which of the two task types
it becomes, and \textsf{KeyClaims} records the paper's terminal quantities used to author the
held-out reference. Papers are mined from high-impact venues and recent work from established
groups across nine scientific domains, giving $5{,}878$ candidates. Filtering reads the extracted
fields: a paper is dropped when its \textsf{Experiment} is purely wet-lab or its data is not
publicly available (the task could not be run), and when its \textsf{Method} and
\textsf{Motivation} describe a single-step lookup rather than an investigation that rewards
multiple stages. From what remains, $N=100$ tasks are authored, drawn from seven of the domains.
Every task carries two labels: its domain, and one of the \emph{novelty-move} types, read from its \textsf{Tension} and \textsf{Conclusion}, recording what kind of
contribution the paper makes. Tasks are restricted to papers from 2024 onward. This applies to the open-ended discovery subset; the target-anchored
optimization subset follows the source benchmark's task selection and includes earlier papers (Appendix~\ref{app:tasks}).


\subsubsection{Task types and composition}
\label{sec:stats}
A task's \textsf{Conclusion} and \textsf{Motivation} also fix its type, according to whether the
paper's goal terminates in a quantity computable on held-out data or in a qualitative finding that
does not. \emph{Open-ended discovery} tasks ($n=70$) have no such quantity: a human reference
exists, but nothing in the environment tells the agent whether it is getting closer, and the space
of acceptable methods is wide. \emph{Target-anchored optimization} tasks ($n=30$) expose an
explicit objective---a human state of the art or a computable metric---giving a well-posed
direction of improvement, though the method to reach it is still unspecified. The split is not free of content: an explicit metric exists only for certain kinds
of contribution, so every target-anchored task is a new-regime, incremental, or method-correction
move, while the reconciliation, mechanism, consensus-overturn, and scaling-relation moves are
open-ended only. This also shows on the domain axis (\Cref{fig:stats}): the open-ended set spans the seven represented domains, whereas the target-anchored set concentrates in
the domains where a computable reference quantity is available. The two types are therefore not
matched on domain or move, and are reported separately throughout---never compared across either.

 
\begin{figure}[t]
    \centering
    \includegraphics[width=0.95\textwidth]{figs/fig_domain_novelty.pdf}
    \caption{Composition of \bench ($N=100$), each axis split by task type: scientific domain
    (left) and novelty-move type (right). Open-ended discovery in blue, target-anchored
    optimization in terracotta.}
    \label{fig:stats}
\end{figure}
\subsection{AutoResearch Trajectories}
\subsubsection{Six-stage end-to-end auto research rollout}
\label{sec:rollout}

Each task is run as one autonomous rollout. Given only the query and a fresh sandbox with code
execution enabled, the agent works without human intervention through six stages---ideation and
planning, retrieval and synthesis, execution and implementation, analysis and interpretation,
writing and documentation, and a final self-verification and review pass over its own
report---so every task involves real retrieval, computation, and artifact generation rather than
answer lookup. The rollout ends when the agent commits a final report or exhausts its budget, and
we log the complete trajectory: the full interaction and tool-call trace, all generated code and
data, intermediate outputs, and the final report.

\subsubsection{Harness-model combinations}
\label{sec:setup:agents}
An agent is a \emph{harness--model} pair: the harness supplies the tool loop, file system, code
execution, and control flow, and the backbone model drives it. We evaluate the frontier harnesses
\textbf{Claude Code}, \textbf{Codex}, and \textbf{Gemini CLI} against backbones from several model
families (\Cref{tab:harness-model}). Running one backbone under several harnesses, and one harness
over several backbones, separates failures of the scaffold from failures of the model and shows
whether a failure mode is idiosyncratic to one system or shared. Because the central analysis is the
failure-pattern taxonomy rather than a ranking this suffices---the failure mechanisms recur across
combinations, and any claim that depends on a particular combination or subset is stated with its
population.

\begin{table}[t]
\centering\small
\caption{Harness--model combinations under evaluation.}
\label{tab:harness-model}
\begin{tabular}{@{}ll@{}}
\toprule
\textbf{Harness} & \textbf{Backbone models} \\
\midrule
Claude Code & opus-4.8, claude-sonnet-5, qwen3.7-max, \\
            & glm-5.2, minimax-m3, deepseek-v4-pro \\
Codex       & gpt-5-mini \\
Gemini CLI  & gemini-3.5-flash \\
\bottomrule
\end{tabular}
\end{table}

\subsubsection{Trajectory statistics}
\label{sec:setup:rollout}
Every task is run by every harness--model combination, with the query as its only input, which
over the $8$ combinations yields $800$ trajectories comprising
$73$k tool calls at an average of $92.3$ steps per episode. The complete trajectory---not the
final report alone---is the unit of the process-level analysis in the sections that follow.
Sandbox configuration, budgets, prompts, and the full task manifest are in
Appendix~\ref{app:extraction}.